\documentclass[10pt, aspectratio=169]{beamer}

% Utilisation d'un thème simple et moderne
\usetheme{Madrid}

\usepackage[utf8]{inputenc}
\usepackage[T1]{fontenc}
\usepackage[french]{babel}
\usepackage{graphicx}
\usepackage{enumitem}

%--- Metadonnées du document
\title{Le Cyberharcèlement à l'Ère de l'IA}
\subtitle{Risques, Enjeux et Protection}
\author{Maël Mignard}
\institute{BTS SIO}
\date{\today}

\begin{document}

\frame{\titlepage}

\begin{frame}{Sommaire}
  \tableofcontents
\end{frame}

%-------------------------------------------------------------------------------
\section{Introduction}
%-------------------------------------------------------------------------------

\begin{frame}{Qu'est-ce que le Cyberharcèlement ?}
  \begin{block}{Définition}
    Le cyberharcèlement est une agression répétée via des supports numériques (réseaux sociaux, messages, etc.).
    \pause
    Sa particularité : il peut toucher une victime n'importe où, n'importe quand.
  \end{block}
  
  \begin{alertblock}{Caractéristiques}
    \begin{itemize}
      \item \textbf{Anonymat :} L'écran peut encourager les agresseurs.
      \item \textbf{Viralité :} Un contenu blessant se propage très vite.
      \item \textbf{Persistance :} Les traces restent en ligne longtemps.
    \end{itemize}
  \end{alertblock}
\end{frame}

%-------------------------------------------------------------------------------
\section{L'IA : Nouveaux Risques}
%-------------------------------------------------------------------------------

\begin{frame}{Comment l'IA aggrave le problème}
  
  \begin{columns}[T]
    \begin{column}{.5\textwidth}
      \begin{block}{Deepfakes \& Usurpation}
        L'IA peut créer de fausses vidéos ou de faux messages audio pour humilier ou manipuler.
        \begin{itemize}
          \item Fausses déclarations.
          \item Revenge porn.
        \end{itemize}
      \end{block}
    \end{column}
    \begin{column}{.5\textwidth}
      \begin{alertblock}{Automatisation de la Haine}
        Des bots (robots logiciels) peuvent être programmés pour envoyer massivement des messages de haine à une victime.
      \end{alertblock}
    \end{column}
  \end{columns}

  \begin{block}{Doxing assisté par IA}
    L'IA peut collecter des informations personnelles en ligne pour identifier et nuire à quelqu'un.
  \end{block}
\end{frame}

%-------------------------------------------------------------------------------
\section{Enjeux Éthiques}
%-------------------------------------------------------------------------------

\begin{frame}{Les Dilemmes de l'IA}
  \begin{alertblock}{Biais des Algorithmes}
    Une IA entraînée avec des données biaisées peut reproduire et amplifier les discriminations (racisme, sexisme...). 
  \end{alertblock}

  \begin{columns}[T]
    \begin{column}{.5\textwidth}
      \begin{block}{La "Boîte Noire"}
        Il est parfois impossible de savoir comment une IA prend une décision. C'est un problème de transparence.
      \end{block}
    \end{column}
    \begin{column}{.5\textwidth}
      \begin{block}{Responsabilité}
        Qui est responsable si une IA cause un dommage ?
        \begin{itemize}
          \item Le développeur ?
          \item L'utilisateur ?
          \item L'IA elle-même ?
        \end{itemize}
      \end{block}
    \end{column}
  \end{columns}
\end{frame}

%-------------------------------------------------------------------------------
\section{Comment se Protéger ?}
%-------------------------------------------------------------------------------

\begin{frame}{Pistes de Solution}

  \begin{columns}[T]
    \begin{column}{.5\textwidth}
      \begin{block}{La Loi}
        \begin{itemize}
          \item \textbf{IA Act (Europe) :} Régule les IA selon leur niveau de risque.
          \item \textbf{Lois sur le cyberharcèlement :} Sanctions plus lourdes.
        \end{itemize}
      \end{block}
    \end{column}
    \begin{column}{.5\textwidth}
      \begin{alertblock}{Notre Rôle de Technicien}
        En tant que futurs BTS SIO, nous devons :
        \begin{itemize}
          \item Créer des systèmes sécurisés.
          \item Promouvoir une IA éthique.
          \item Sensibiliser les utilisateurs.
        \end{itemize}
      \end{alertblock}
    \end{column}
  \end{columns}

  \begin{block}{Éducation et Esprit Critique}
    La meilleure défense est la prévention : éduquer aux bons usages du numérique et développer son esprit critique.
  \end{block}
\end{frame}

%-------------------------------------------------------------------------------
\section{Conclusion}
%-------------------------------------------------------------------------------

\begin{frame}{Conclusion}
  L'IA offre de nouvelles possibilités, mais crée aussi des risques pour le harcèlement en ligne.
  \vfill
  Il faut une approche qui combine :
  \begin{itemize}
      \item \textbf{Lois} (régulation)
      \item \textbf{Technologie responsable}
      \item \textbf{Éducation}
  \end{itemize}
  \vfill
  En tant que techniciens, nous sommes en première ligne.
\end{frame}

%-------------------------------------------------------------------------------
\begin{frame}
    \centering
    \Huge\bfseries Questions ?
\end{frame}
%-------------------------------------------------------------------------------

\end{document}
